\chapter{Esercizi: Ricerca e ordinamento in Java}
\label{chap:esercizi-ricerca-ordinamento-java}

\begin{exercise}
\label{ex:fi-esercizi-java-alg-linear-prima}
Che cosa restituisce \texttt{linearSearch()} se il valore compare
più volte o non compare nell'array?
\end{exercise}

Restituisce l'indice della prima occorrenza, oppure \(-1\) se il
valore è assente (\cref{sec:fi-java-alg-lineare}).

\begin{exercise}
\label{ex:fi-esercizi-java-alg-linear-costi}
Quali sono la complessità temporale migliore e peggiore e lo
spazio ausiliario della ricerca lineare?
\end{exercise}

Per \(n\ge1\), il caso migliore richiede \(\Theta(1)\), il peggiore
\(\Theta(n)\). Lo spazio ausiliario è \(\Theta(1)\)
(\cref{sec:fi-java-alg-lineare}).

\begin{exercise}
\label{ex:fi-esercizi-java-alg-binaria-precondizione}
Quale precondizione deve rispettare l'array passato ai metodi
di ricerca binaria?
\end{exercise}

L'array deve essere ordinato in senso non decrescente; altrimenti
l'esclusione di metà dell'intervallo può eliminare il valore
cercato (\cref{sec:fi-java-alg-binaria}).

\begin{exercise}
\label{ex:fi-esercizi-java-alg-binaria-spazio}
Perché le due versioni di ricerca binaria hanno lo stesso costo
temporale asintotico, ma un diverso costo spaziale?
\end{exercise}

Entrambe dimezzano l'intervallo a ogni passo, richiedendo
\(O(\log n)\) tempo nel caso peggiore. La versione iterativa
conserva soltanto gli indici, quindi usa \(\Theta(1)\) spazio;
quella ricorsiva richiede una pila di chiamate di profondità
\(O(\log n)\)
(\cref{sec:fi-java-alg-binaria-iterativa,sec:fi-java-alg-binaria-ricorsiva}).

\begin{exercise}
\label{ex:fi-esercizi-java-alg-binaria-duplicati}
La ricerca binaria restituisce sempre la prima occorrenza di un
valore ripetuto?
\end{exercise}

No. Restituisce un indice in cui compare il valore cercato,
non necessariamente quello della prima occorrenza
(\cref{sec:fi-java-alg-binaria}).

\begin{exercise}
\label{ex:fi-esercizi-java-alg-selezione-supporto}
A che cosa servono i metodi \texttt{findMinPos()} e \texttt{swap()}
nell'ordinamento per selezione?
\end{exercise}

Il primo individua la posizione del minimo nel tratto non
ordinato; il secondo scambia tale elemento con quello iniziale
del tratto (\cref{sec:fi-java-alg-selezione}).

\begin{exercise}
\label{ex:fi-esercizi-java-alg-selezione-confronti}
Perché \texttt{selectionSort()} richiede tempo quadratico anche
quando l'array è già ordinato?
\end{exercise}

Continua a cercare il minimo in ogni suffisso e svolge
\(n(n-1)/2\) confronti, indipendentemente dall'ordine iniziale
(\cref{sec:fi-java-alg-selezione}).

\begin{exercise}
\label{ex:fi-esercizi-java-alg-inserimento-invariante}
Quale proprietà mantiene \texttt{insertionSort()} all'inizio
di ogni iterazione esterna?
\end{exercise}

Il prefisso che precede l'elemento corrente è già ordinato.
L'elemento viene inserito nella posizione corretta conservando
questa proprietà (\cref{sec:fi-java-alg-inserimento}).

\begin{exercise}
\label{ex:fi-esercizi-java-alg-inserimento-casi}
Quando \texttt{insertionSort()} ha costo lineare e quando
quadratico?
\end{exercise}

Se l'array è ordinato, non servono spostamenti e il tempo è
\(\Theta(n)\). In ordine strettamente decrescente occorrono
\(n(n-1)/2\) spostamenti, quindi \(\Theta(n^2)\)
(\cref{sec:fi-java-alg-inserimento}).

\begin{exercise}
\label{ex:fi-esercizi-java-alg-inserimento-stabile}
Perché l'ordinamento per inserimento è stabile?
\end{exercise}

Il ciclo interno sposta soltanto elementi strettamente maggiori
di quello da inserire; gli elementi uguali conservano l'ordine
relativo (\cref{sec:fi-java-alg-inserimento}).

\begin{exercise}
\label{ex:fi-esercizi-java-alg-merge-caso-base}
Qual è il caso base di \texttt{mergeSort()} e perché garantisce
la terminazione?
\end{exercise}

Un array di lunghezza zero o uno è già ordinato. Ogni invocazione
ricorsiva riceve una metà strettamente più piccola, perciò
raggiunge il caso base (\cref{sec:fi-java-alg-fusione}).

\begin{exercise}
\label{ex:fi-esercizi-java-alg-merge-residui}
Perché \texttt{merge()} copia separatamente gli elementi
rimanenti dopo il primo ciclo?
\end{exercise}

Il primo ciclo termina quando una delle due metà ordinate è
esaurita. Gli elementi residui dell'altra metà sono già ordinati
e non precedono quelli copiati: possono quindi essere trasferiti
direttamente (\cref{sec:fi-java-alg-fusione}).

\begin{exercise}
\label{ex:fi-esercizi-java-alg-merge-costi}
Quali sono i costi temporale e spaziale dell'implementazione
Java di MergeSort?
\end{exercise}

Il tempo è \(\Theta(n\log n)\) in tutti i casi; le copie
temporanee richiedono \(\Theta(n)\) spazio ausiliario
complessivo al picco (\cref{sec:fi-java-alg-fusione}).

\begin{exercise}
\label{ex:fi-esercizi-java-alg-quick-partizione}
Che cosa garantisce \texttt{partition()} quando restituisce
l'indice del pivot?
\end{exercise}

Il pivot si trova nella posizione definitiva; tutti gli elementi
a sinistra non lo superano e quelli a destra sono maggiori
(\cref{sec:fi-java-alg-rapido}).

\begin{exercise}
\label{ex:fi-esercizi-java-alg-quick-peggiore}
Perché scegliere l'ultimo elemento come pivot può produrre
un costo quadratico per QuickSort?
\end{exercise}

Con input già ordinato, la partizione può produrre un intervallo
vuoto e uno di lunghezza \(n-1\). Ripetendo tale squilibrio,
la somma dei costi lineari delle partizioni è \(\Theta(n^2)\)
(\cref{sec:fi-java-alg-rapido}).

\begin{exercise}
\label{ex:fi-esercizi-java-alg-quick-spazio}
Perché lo spazio ausiliario di QuickSort dipende dalle
partizioni, pur effettuando gli scambi nello stesso array?
\end{exercise}

La partizione usa \(\Theta(1)\) spazio, ma le chiamate ricorsive
rimangono sulla pila. La profondità è \(O(\log n)\) con
partizioni bilanciate e può raggiungere \(O(n)\) nel caso
peggiore (\cref{sec:fi-java-alg-rapido}).

\begin{exercise}
\label{ex:fi-esercizi-java-alg-stabilita}
Quali algoritmi tra i quattro ordinamenti presentati sono
stabili nelle rispettive implementazioni?
\end{exercise}

Insertion Sort e MergeSort. Selection Sort e QuickSort non
garantiscono la stabilità (\cref{chap:ricerca-ordinamento-java}).
